\section{Discussion}
\label{sec:discussion}

\subsection{Designated and Plural Authority}
\label{sec:thesis}

This paper covers one class of decision: a command decision, which gets its legitimacy
from one designated person's role, takes effect when it is issued, and is settled by
authority, never by merge, when it conflicts with another. It binds when it is admitted to
the authoritative journal under a valid grant. Decisions whose authority is shared or
contested belong to the opposite pole, where legitimacy comes from agreement among
several parties and a decision binds only at reconciliation; there, merge or consensus is
the right tool. The class of decision, not the state of the network, chooses between the
two. Our earlier design let reachability choose the commander; the model of
\S~\ref{sec:system-model} removes that.

The boundary is also where this paper stops. \systemname{} defines residual decisions,
attributes them and labels them; it does not resolve them. Deciding what happens to an
order issued under a superseded grant---confirm it, reverse it, or let the new commander
rule on it---is a reconciliation that acts on command decisions, and we leave it to
future work that composes the two poles.

\subsection{Security Assumptions}

The target deployment assumes an organizational CA with enrolled devices and users, and
the model assumes credentials are not stolen. Revocation of a stolen tablet or vehicle
node is eventual under total disconnection, and a captured command edge can poison the
journal for as long as its holder's grant is valid; a quorum-inferred design fails
differently, not better, by moving authority away from the designated commander. The
design relies on signed grants, fencing and audit, not on Byzantine tolerance.

\subsection{Evaluation Boundaries and Non-claims}
\label{sec:non-claims}
\label{sec:alternatives}

The paper does not claim that \systemname{} solves succession under partition: the
trade-off rules that out, and the prototype realizes only S1. The claims about S0 and
S2--S5 in Table~\ref{tab:policies} are hypotheses until the comparison of
\S~\ref{sec:eval-policies} has been run.

Both formal models cover coordinated handover only. The TLA+ model has one shared
journal and promotes only when the old command edge is reachable; the Hypothesis handoff
model promotes only after the old command edge has forwarded its whole journal; neither
has a lease. Model checking therefore establishes the six invariants for that case. The
cut-off case of Figure~\ref{fig:promotion-timeline} has a runtime witness only, and in it
\textit{DurabilityAcrossPromotion} does not hold for entries the old command edge never
forwarded.

The prototype's gaps against the model are listed in \S~\ref{sec:model-gap}. In
addition, \systemname{} records decisions but does not deliver them to responder edges;
any delivery acknowledgement must not block, since a two- or three-phase commit would
stall every decision while a team is out of coverage (R1). Semantic rejection of
conflicting values is not implemented.

The emulation isolates the authority and replication mechanisms from deployment
variables: physical mesh transport, security overlay cost, and Android client
integration. It does not claim physical Rajant mesh behavior, measured WireGuard or mTLS
overhead, or end-to-end Android validation, and the workload is single-incident and
limited-scale. The connectivity schedule is not a mobility model: no vehicle positions,
path loss, or RAN topology are simulated. Command-side writes are exercised through
\textit{syncd}'s accept path rather than the command \textit{ops-api}, and core bootstrap
and backfill are synthetic rather than measured.
